% ===================================================================
%  TABELO N° 3 — Puereso e Yuneso.
%  Feuillets 19 a 25 du fac-simile = folios 17 a 23.
%  Correspondance : folio imprime = numero de feuillet - 2.
%
%  Ca tabelo (kom la 4-ma) kombinas « ceni familiala », quale la
%  parentezo che l'apertajo l'expresas. Ol propra kontenas du ceni —
%  « Unesma ceno » (Bapto-ceremonio en vilajo) e « Duesma ceno »
%  (Festo publika) — kada un kun sua propra numerado d'alineas,
%  ripetanta de 1. La cles %%K sequas la numerado quale imprimita :
%  t03-01-1 ... t03-16-1 aparas do dufoye, unfoye por singla ceno —
%  lo esas l'imprimo mem qua restartigas, ne la transskripto (videz
%  texto/io/17-tabelo-08.tex e 18-tabelo-09.tex por la sama padrono).
% ===================================================================

% --- Feuillet 19 (folio 17 : apertajo di la tabelo) -----------------
\begin{VUpage}[19]{}
%%K t03-apar-1 apar
\VUsaut{3.0mm}
\VUtitre{15.0pt}{60}{TABELO N\textsuperscript{o} 3}
\VUsaut{1.6mm}
\VUtitre{11.4pt}{0}{Puereso e Yuneso.}
\VUsaut{3.4mm}

%%K t03-apar-2 apar
\VUblancAlinea
(La 3 - ma e 4 - ma tabeli esas tale kom\cc
binita ke li prizentas en quar \VUgras{ceni familiala}\nl
l'esencal nocioni pri la \VUgras{vetero}, l'\VUgras{evo}, la \VUgras{fami}\cc
\VUgras{lio}, la \VUgras{vestaro} e la \VUgras{stando.})

\VUsaut{4.0mm}
%%K t03-tit-1 sub
\VUcentre{12.0pt}{0}{\textit{Unesma ceno.}}
\VUsaut{3.0mm}

%%K t03-tit-2 sub
\VUpk{10.2pt}{40}{Bapto-ceremonio en vilajo.}
\VUsaut{2.4mm}

%%K t03-tit-3 sub
\VUpk{10.2pt}{40}{La printempo.}
\VUsaut{3.0mm}

%%K t03-c1-01-1 p
\VUblancAlinea
1. --- Ni esas en la \VUgras{printempo}, en la monati\nl
\VUgras{marto}, \VUgras{aprilo} o \VUgras{mayo}, en la dii di la \VUgras{semano},\nl
\VUgras{lundio}, \VUgras{mardio}, o \VUgras{merkurdio}. Esas \VUgras{dek kloki}\nl
matine.

%%K t03-c1-02-1 p
\VUblancAlinea
2. --- La \VUgras{vetero} esas bela, l'\VUgras{aero} pura. La\nl
suno brilas. Kelka \VUgras{nubeti} \textsuperscript{(2)} acensas en la\nl
blua \VUgras{cielo} \textsuperscript{(1)}.

%%K t03-c1-03-1 p
\VUblancAlinea
3. --- La \VUgras{arbori} apene havas \VUgras{folii}. La tenua\nl
\VUgras{popli} \textsuperscript{(3)} e la alta \VUgras{platano} \textsuperscript{(17)} til nun havas nur\nl
burjoni.

%%K t03-c1-04-1 p
\VUblancAlinea
4. --- La \VUgras{hirundi} \textsuperscript{(4)} rivenas e flugetas kun\nl
joyoza pipiado cirkum la \VUgras{flecho} \textsuperscript{(7)} dil \VUgras{klo}\cc
\VUgras{sheyo} \textsuperscript{(5)} qua kronizas la \VUgras{kirko} \textsuperscript{(6)} dil vilajo.

%%K t03-tit-4 sub
\VUsaut{3.4mm}
\VUpk{10.2pt}{40}{La Kirko.}
\VUsaut{3.0mm}

%%K t03-c1-05-1 p
\VUblancAlinea
5. --- Tre simpla esas ta kirko kun lua\nl
granda \VUgras{portalo} \textsuperscript{(13)} e grosa \VUgras{horlojo} \textsuperscript{(8)}. La mi\cc
kra \VUgras{indexo} \textsuperscript{(9)} esas sur la cifro X, ol indikas la
\parplein
\end{VUpage}

% --- Feuillet 20 (folio 18) ------------------------------------------
\begin{VUpage}[20]{18}
%%K t03-c1-05-1 p suite
\VUcontinue
\VUgras{kloki.} La \VUgras{granda} \textsuperscript{(10)} esas sur XII (dimezo); ol\nl
indikas la \VUgras{minuti}.

%%K t03-c1-06-1 p
\VUblancAlinea
6. --- En la klosheyo, la \VUgras{kloshi} \textsuperscript{(11)} gaye sona\cc
das. Ye la somito dil tekto stacas mikra \VUgras{kru}\cc
\VUgras{co} \textsuperscript{(12)} petra.

%%K t03-c1-07-1 p
\VUblancAlinea
7. --- La kirko havas ne nur granda \VUgras{por}\cc
\VUgras{talo} \textsuperscript{(13)}, ol havas anke pordeto laterala. Tra\nl
ica pordeto sioro \VUgras{paroko} \textsuperscript{(15)} vestizita per sua\nl
\VUgras{sutano}, ekiras la \VUgras{sakristio} \textsuperscript{(14)} e adiras la \VUgras{paro}\cc
\VUgras{keyo} \textsuperscript{(19)}.

%%K t03-c1-08-1 p
\VUblancAlinea
8. --- Proxim ilu, yen la \VUgras{laktovendistino} \textsuperscript{(24)}\nl
kun sua \VUgras{asno} \textsuperscript{(23)}. El retrovenas de la urbo,\nl
en qua el kunportis bona lakto por l'infanti.

%%K t03-tit-5 sub
\VUsaut{4.0mm}
\VUpk{10.2pt}{40}{La frukto-gardeno.}
\VUsaut{3.4mm}

%%K t03-c1-09-1 p
\VUblancAlinea
9. --- Sinioro Aliboron (l'asno) esas tote kon\cc
tenta pri ta matinal kuro. Lu delicoze aspiras\nl
l'aero balzamita da la parfumo dil \VUgras{flori} \textsuperscript{(20)} qui\nl
kovras la \VUgras{fruktarbori} \textsuperscript{(18)} dil vicina \VUgras{frukto-gar}\cc
\VUgras{deno}. Tote rozea e blanka esas ta \VUgras{persikieri} \textsuperscript{(18)}\nl
ed \VUgras{abrikotieri} \textsuperscript{(19)}, e li esperigas bona rekolto.

%%K t03-c1-10-1 p
\VUblancAlinea
10. --- Dume, la mikra \VUgras{abeli} \textsuperscript{(22)} dil \VUgras{abe}\cc
\VUgras{luyo} \textsuperscript{(21)} iras ibe marode kolektar sua dolca\nl
\VUgras{mielo} zumante.

%%K t03-c1-11-1 p
\VUblancAlinea
11. --- Ma quo eventas ? Yen ke la pueri dil\nl
vilajo adkuras e fugigas la \VUgras{gansi} \textsuperscript{(25)} pavori\cc
gita. To esas la ekiro del kirko da la sequan\cc
taro pos la \VUgras{bapteso} dil filio di la notario.

%%K t03-c1-12-1 p
\VUblancAlinea
12. --- La mikra Iakobus, en sua \VUgras{bersilo} \textsuperscript{(26)},\nl
vekas pro la joyoza sonigado dil kloshi, e lua\nl
fratino Magdalena qua, el anke, volas vidar la\nl
bapto-sequantaro, pregas sua matro venar por
\parplein
\end{VUpage}

% --- Feuillet 21 (folio 19) ------------------------------------------
\begin{VUpage}[21]{19}
%%K t03-c1-12-1 p suite
\VUcontinue
pektar elua hararo e metigar elua vestaro ye\nl
la solio dil domo.

%%K t03-c1-13-1 p
\VUblancAlinea
13. --- La matro asentas. El metas sua \VUgras{fu}\cc
\VUgras{lardo} \textsuperscript{(28)}, sua \VUgras{korsajo} \textsuperscript{(29)} e sua \VUgras{avantalo} \textsuperscript{(30)},\nl
ed el venas sidar sur mikra stulo avan la pordo.\nl
Magdalena havas nur sua \VUgras{subjupo} \textsuperscript{(31)} e sua\nl
\VUgras{korseto} \textsuperscript{(32)}.

%%K t03-c1-14-1 p
\VUblancAlinea
14. --- Quante felica el esas ! El oblivias sua\nl
favorata flori, sua \VUgras{dianti} \textsuperscript{(34)}, sur qui pozas su\nl
la \VUgras{papilioni} \textsuperscript{(35)}, sua \VUgras{tulipi} \textsuperscript{(36)}, sua \VUgras{maiflori} \textsuperscript{(37)},\nl
en \VUgras{poti} \textsuperscript{(38)}, sur la \VUgras{muro} \textsuperscript{(33)}, e sua \VUgras{rozi} \textsuperscript{(39)} qui\nl
formacas hego tante bela. El konsideras la\nl
richa vestaro di la dami dil eskorto.

%%K t03-c1-15-1 p
\VUblancAlinea
15. --- La pueri dil vilajo ne tre atencas la\nl
vestaro. Li precipitas su e shokpulsas l'una\nl
l'altra por havar \VUgras{bonboni} \textsuperscript{(55)} o \VUgras{monetpeci} \textsuperscript{(52)}.\nl
Un de li ploras \textsuperscript{(40)}.

%%K t03-c1-16-1 p
\VUblancAlinea
16. --- L'altra marchis sur la pedo di la\nl
\VUgras{hundo} \textsuperscript{(42)} qua kriante fugas e fugigas la \VUgras{ha}\cc
\VUgras{nino} \textsuperscript{(43)} e lua \VUgras{hanyuneti} \textsuperscript{(44)}.

%%K t03-tit-6 sub
\VUsaut{5.0mm}
\VUpk{10.2pt}{40}{La baptal defilo.}
\VUsaut{4.0mm}

%%K t03-c1-17-1 p
\VUblancAlinea
17. --- Avan l'eskorto marchas la \VUgras{nutris}\cc
\VUgras{tino} \textsuperscript{(45)}. El havas bela \VUgras{kofio} \textsuperscript{(46)} kun longa ru\cc
bandi. El portas la \VUgras{novenaskinto} \textsuperscript{(47)} tote ves\cc
tizita per \VUgras{denteli} \textsuperscript{(48)}.

%%K t03-c1-18-1 p
\VUblancAlinea
18. --- Dop la nutristino venas la \VUgras{baptopa}\cc
\VUgras{trulo} \textsuperscript{(49)} e la \VUgras{baptomatro} \textsuperscript{(53)}. Li disdonas la\nl
bonboni e monetpeci. La baptopatrulo havas\nl
\VUgras{ganti} \textsuperscript{(51)} e \VUgras{cilindra chapelo} \textsuperscript{(50)}. La baptomatro\nl
manue tenas bela \VUgras{buketo} \textsuperscript{(54)}.
\end{VUpage}

% --- Feuillet 22 (folio 20) ------------------------------------------
\begin{VUpage}[22]{20}
%%K t03-c1-19-1 p
\VUblancAlinea
19. --- Dop li, ni vidas la \VUgras{onklulo} \textsuperscript{(56)} e la\nl
\VUgras{onklino} \textsuperscript{(58)} dil infanto. L'onklino havas ele\cc
ganta \VUgras{chapelo} \textsuperscript{(59)}, l'onklulo nigra \VUgras{redingoto} \textsuperscript{(57)}\nl
e blanka jileto. Yen pose la \VUgras{kuzulo} \textsuperscript{(60)} e la\nl
\VUgras{kuzino} \textsuperscript{(61)}. Elca tenas per la manuo la \VUgras{fra}\cc
\VUgras{tino} \textsuperscript{(62)} dil novenaskinto, la mikra Berta.

%%K t03-c1-20-1 p
\VUblancAlinea
20. --- L'altra \VUgras{fratulo} \textsuperscript{(63)} di Berta marchas\nl
dope. Il donas la manuo a \VUgras{parentino} \textsuperscript{(64)}. La\nl
\VUgras{amiki} \textsuperscript{(65)} di la familio venas pose.

%%K t03-tit-7 sub
\VUsaut{4.6mm}
\VUpk{10.2pt}{40}{La spektanti.}
\VUsaut{3.4mm}

%%K t03-c1-21-1 p
\VUblancAlinea
21. --- Proxim la sequantaro, yen \VUgras{mendi}\cc
\VUgras{kero} \textsuperscript{(66)} apogata sur sua \VUgras{grucho} \textsuperscript{(67)}. Il deman\cc
das almono (mendikas).

%%K t03-c1-22-1 p
\VUblancAlinea
22. --- Altralatere esas mikra puerulo tre\nl
\VUgras{polita} \textsuperscript{(68)}. Il salutas e deziras « \VUgras{bon jorno} » al\nl
personi quin il konocas.

%%K t03-c1-23-1 p
\VUblancAlinea
23. --- La vilajani ekiris sua hemi por vidar\nl
la baptal defilo. Ni vidas \VUgras{rurano} \textsuperscript{(69)} kun sua\nl
\VUgras{kotona boneto} ed apude \VUgras{ruranino} \textsuperscript{(70)}. Pose \VUgras{ol}\cc
\VUgras{dino} \textsuperscript{(71)} apogata sur sua \VUgras{bastono} \textsuperscript{(72)}. Fine la\nl
\VUgras{muelisto} \textsuperscript{(73)} kun sua granda \VUgras{felta chapelo} \textsuperscript{(74)}\nl
e yuna ruranino pedvestizita per \VUgras{ligna shui} \textsuperscript{(75)},\nl
qua docas marchar a sua infanteto.

%%K t03-c1-24-1 p
\VUblancAlinea
24. --- Ta ceno esas tre amuzanta.

\VUsaut{4.0mm}
%%K t03-tit-8 sub
\VUcentre{12.0pt}{0}{\textit{Duesma ceno.}}
\VUsaut{3.0mm}

%%K t03-tit-9 sub
\VUpk{10.2pt}{40}{Festo publika.}
\VUsaut{2.4mm}

%%K t03-tit-10 sub
\VUpk{10.2pt}{40}{Navigal ludi e regati.}
\VUsaut{3.0mm}

%%K t03-c2-01-1 p
\VUblancAlinea
1. --- La \VUgras{somero} eventis. La varmega suno\nl
dil monato \VUgras{junio} (julio od \VUgras{agosto}) instigas la\nl
yuni a balnar e kanotagar.
\end{VUpage}

% --- Feuillet 23 (folio 21) ------------------------------------------
\begin{VUpage}[23]{21}
%%K t03-c2-02-1 p
\VUblancAlinea
2. --- Sur la dextra rivo dil fluvio, la kano\cc
teri desvestizas su por partoprenar en la navi\cc
gal ludi e regati.

%%K t03-c2-03-1 p
\VUblancAlinea
3. --- Un de li desmetas sua \VUgras{shui} \textsuperscript{(1)}; il des\cc
ligas (desbutonagas) oli. Lua du kamaradi ja\nl
esas pronta. De nun li havas nur sua \VUgras{trikota}\nl
\VUgras{vesto} \textsuperscript{(2)} e \VUgras{balnokalsono} \textsuperscript{(3)}. Li konversas var\cc
tante sua amiki. Li parolas pri la ferio e la\nl
festo qua eventas sur l'altra rivo.

%%K t03-c2-04-1 p
\VUblancAlinea
4. --- An la sulo, proxim li, jacas la \VUgras{vesti} di\nl
lia kamaradi : \VUgras{paro} de \VUgras{boteti} \textsuperscript{(4)}, \VUgras{shui} \textsuperscript{(5)}, \VUgras{kal}\cc
\VUgras{zeti} \textsuperscript{(6)}, \VUgras{felta} \textsuperscript{(7)} e \VUgras{palia} \textsuperscript{(8)} chapeli e \VUgras{kravato} \textsuperscript{(9)}.

%%K t03-c2-05-1 p
\VUblancAlinea
5. --- Un de la yunuli sorgoze faldis sua\nl
\VUgras{vestio} \textsuperscript{(10)}. Il pozas olu sur tualet-tuko, en qua\nl
lua vicino balde inkluzos la du \VUgras{vesti}.

%%K t03-c2-06-1 p
\VUblancAlinea
6. --- Sisesma puero tardesas. Il havas an\cc
kore sua \VUgras{kasqueto} \textsuperscript{(11)} sur la kapo. Il desmetis\nl
nek sua \VUgras{kamizo} \textsuperscript{(12)}, nek sua \VUgras{kolumo} \textsuperscript{(13)}, nek\nl
sua \VUgras{mansheti} \textsuperscript{(14)}. Il manue tenas sua \VUgras{jileto} \textsuperscript{(17)}\nl
e havas ankore sua \VUgras{pantalono} \textsuperscript{(16)} e \VUgras{breteli} \textsuperscript{(15)}.\nl
Certe il ne esos pronta por la kuro proxima.

%%K t03-c2-07-1 p
\VUblancAlinea
7. --- Proxim la yunuli, voyeto sur la lateri\nl
di qua esas multa \VUgras{kardoni} \textsuperscript{(18)}, duktas al bordo\nl
di la rivero, ube la patrulo Iakobus preparas\nl
inter la \VUgras{kani} \textsuperscript{(19)} la \VUgras{pirotekno} \textsuperscript{(20)} quan on pafos\nl
vespere.

%%K t03-tit-11 sub
\VUsaut{4.4mm}
\VUpk{10.2pt}{40}{La kuro.}
\VUsaut{3.4mm}

%%K t03-c2-08-1 p
\VUblancAlinea
8. --- Sur la rivero, kelka dami, shirmanta\nl
su per sua parasuni, promenas per \VUgras{barko} \textsuperscript{(21)}.\nl
Eli sequas la kuro qua esas tre interesanta. En\nl
singla \VUgras{yolo} \textsuperscript{(22)}, quar \VUgras{remeri} \textsuperscript{(24)} remas per sua
\parplein
\end{VUpage}

% --- Feuillet 24 (folio 22) ------------------------------------------
\begin{VUpage}[24]{22}
%%K t03-c2-08-1 p suite
\VUcontinue
tota forci, dum ke la \VUgras{guvernanto} \textsuperscript{(26)} tenas la\nl
\VUgras{guvernilo} \textsuperscript{(25)} e direktas la frajil rembatelo. La\nl
\VUgras{remili} \textsuperscript{(23)} kadence frapas l'aquo e la lejera ba\cc
teli navigas kun rapideso vertijiganta.

%%K t03-c2-09-1 p
\VUblancAlinea
9. --- Kelka yunuli interluktas, proxim la\nl
pilastro dil ponto, pri la premio dil \VUgras{busprito}\cc
\VUgras{masto} \textsuperscript{(28)}. Un de li prudente marchas sur la\nl
flexebla e glitiganta masto por atingar la mi\cc
kra \VUgras{flago} \textsuperscript{(29)} fixigita ye l'extremajo. Lua mal\cc
chançoza \VUgras{konkursanto} \textsuperscript{(30)} falis en l'aquo. Il\nl
natas por tervenar.

%%K t03-c2-10-1 p
\VUblancAlinea
10. --- Plu evoza yunuli \VUgras{(batel)-turniras} \textsuperscript{(32)}\nl
proxim la \VUgras{kayo} \textsuperscript{(33)}. Ta omna ludin asistas\nl
grandanombra \VUgras{publiko} \textsuperscript{(31)} apogata an la \VUgras{para}\cc
\VUgras{peto} dil \VUgras{ponto} \textsuperscript{(27)} ed amasigita sur la \VUgras{rivo}.\nl
Kelka spektanti, tamen, riturnas su por sequar\nl
la ludo dil \VUgras{premio-masto} \textsuperscript{(45)} o por admirar\nl
l'\VUgras{urbestral eskorto} qua venas por la \VUgras{disdono}\nl
dil premii.

%%K t03-c2-11-1 p
\VUblancAlinea
11. --- Unesme, yen kelka \VUgras{bubi} \textsuperscript{(35)} preiranta\nl
la \VUgras{muzikisti} \textsuperscript{(36)}; pose venas la \VUgras{urbestro} \textsuperscript{(37)}, la\nl
adyunti e la \VUgras{komonal konsilantaro} di la urbeto.\nl
Fine marchas la \VUgras{pumpisti} \textsuperscript{(38)}.

%%K t03-tit-12 sub
\VUsaut{5.0mm}
\VUpk{10.2pt}{40}{La ferio.}
\VUsaut{3.6mm}

%%K t03-c2-12-1 p
\VUblancAlinea
12. --- Grandanombra homi esas sur la \VUgras{ferio}\cc
\VUgras{placo} \textsuperscript{(39)}. On venas vidar la diversa \VUgras{baraki} :\nl
la ligna \VUgras{kavali} \textsuperscript{(40)}, la \VUgras{cirko} \textsuperscript{(41)} che qua ja ko\cc
mencis la reprezento; la \VUgras{publika balo} \textsuperscript{(42)}, ube\nl
la yunuli e la yunini dil urbo juas la plezuro\nl
dil \VUgras{dansado}.

%%K t03-c2-13-1 p
\VUblancAlinea
13. --- En la fundo, ni vidas la \VUgras{areni} \textsuperscript{(44)} che\nl
qui la brava \VUgras{matadoro} Hispana luktas kontre
\parplein
\end{VUpage}

% --- Feuillet 25 (folio 23) ------------------------------------------
\begin{VUpage}[25]{23}
%%K t03-c2-13-1 p suite
\VUcontinue
la furioza \VUgras{tauro} \textsuperscript{(45)}, pose la \VUgras{relgliteyo} \textsuperscript{(49)}, la\nl
\VUgras{pafeyo} \textsuperscript{(46)} ube on exercas su uzar la \VUgras{pafarmi}\nl
e pafar al \VUgras{skopo}.

%%K t03-c2-14-1 p
\VUblancAlinea
14. --- Proxime, yen la \VUgras{ferial teatro} \textsuperscript{(47)} en\nl
qua on pleas la \VUgras{komedio} e la \VUgras{dramato}. Tre\nl
proxime, esas la barako dil \VUgras{giganto} \textsuperscript{(48)} e dil\nl
\VUgras{nano.} La \VUgras{staturo} di l'unesma esas du metri plus\nl
dek-e-kin centimetri; la duesma esas apene\nl
tam granda kam infanto.

%%K t03-c2-15-1 p
\VUblancAlinea
15. --- Fine, yen la granda \VUgras{menajerio} \textsuperscript{(50)}\nl
kun lua \VUgras{sovaja bestii}, lua \VUgras{tigrulo} \textsuperscript{(51)} kun po\cc
tenta \VUgras{ungli}, lua \VUgras{leonulo} \textsuperscript{(52)} kun densa \VUgras{kol}\cc
\VUgras{krinaro}, lua du \VUgras{jirafi} \textsuperscript{(53)} e lua \VUgras{elefanto} \textsuperscript{(54)} qua\nl
tante habile uzas sua \VUgras{rostro}.

%%K t03-c2-16-1 p
\VUblancAlinea
16. --- La \VUgras{domtisto} \textsuperscript{(55)} qua dresas ta bestii\nl
esas ye la pordo di la barako.

\VUsaut{6.0mm}
\VUfilet{22mm}
\end{VUpage}
